\documentclass[10pt,a4paper]{amsart}
\usepackage{amsmath,amssymb,amsthm,graphicx}
% Retained minimal preamble for the Stacks Project chapter more-algebra.tex (More on Algebra).
% The chapter file pinned-original-source.tex carries only \input{preamble} on its first line, so no
% preamble text lives inside the chapter itself.  The Stacks Project uses one single shared preamble for
% every chapter; the snapshot stages it identically as sources/stacks-chow/preamble.tex and
% sources/stacks-etale-cohomology/preamble.tex.  The block below is copied mechanically, line for line,
% from sources/stacks-chow/preamble.tex lines 14-18 (the Xy-pic setup) and 154-217 (the native theorem
% declarations together with the mathematical macros).  Nothing was added, removed or redefined; the
% wrapper class and the amsmath/amssymb/amsthm/graphicx base packages are supplied by the assembler.
% Chapter: more-algebra.tex -- More on Algebra.
% For commutative diagrams we use Xy-pic
\usepackage[all]{xy}

% We use 2cell for 2-commutative diagrams.
\xyoption{2cell}

% Theorem environments.
%
\theoremstyle{plain}
\newtheorem{theorem}[subsection]{Theorem}
\newtheorem{proposition}[subsection]{Proposition}
\newtheorem{lemma}[subsection]{Lemma}

\theoremstyle{definition}
\newtheorem{definition}[subsection]{Definition}
\newtheorem{example}[subsection]{Example}
\newtheorem{exercise}[subsection]{Exercise}
\newtheorem{situation}[subsection]{Situation}

\theoremstyle{remark}
\newtheorem{remark}[subsection]{Remark}
\newtheorem{remarks}[subsection]{Remarks}

\numberwithin{equation}{subsection}

% Macros
%
\def\lim{\mathop{\mathrm{lim}}\nolimits}
\def\colim{\mathop{\mathrm{colim}}\nolimits}
\def\Spec{\mathop{\mathrm{Spec}}}
\def\Hom{\mathop{\mathrm{Hom}}\nolimits}
\def\Ext{\mathop{\mathrm{Ext}}\nolimits}
\def\SheafHom{\mathop{\mathcal{H}\!\mathit{om}}\nolimits}
\def\SheafExt{\mathop{\mathcal{E}\!\mathit{xt}}\nolimits}
\def\Sch{\mathit{Sch}}
\def\Mor{\mathop{\mathrm{Mor}}\nolimits}
\def\Ob{\mathop{\mathrm{Ob}}\nolimits}
\def\Sh{\mathop{\mathit{Sh}}\nolimits}
\def\NL{\mathop{N\!L}\nolimits}
\def\CH{\mathop{\mathrm{CH}}\nolimits}
\def\proetale{{pro\text{-}\acute{e}tale}}
\def\etale{{\acute{e}tale}}
\def\QCoh{\mathit{QCoh}}
\def\Ker{\mathop{\mathrm{Ker}}}
\def\Im{\mathop{\mathrm{Im}}}
\def\Coker{\mathop{\mathrm{Coker}}}
\def\Coim{\mathop{\mathrm{Coim}}}

% Boxtimes
%
\DeclareMathSymbol{\boxtimes}{\mathbin}{AMSa}{"02}

%
% Macros for moduli stacks/spaces
%
\def\QCohstack{\mathcal{QC}\!\mathit{oh}}
\def\Cohstack{\mathcal{C}\!\mathit{oh}}
\def\Spacesstack{\mathcal{S}\!\mathit{paces}}
\def\Quotfunctor{\mathrm{Quot}}
\def\Hilbfunctor{\mathrm{Hilb}}
\def\Curvesstack{\mathcal{C}\!\mathit{urves}}
\def\Polarizedstack{\mathcal{P}\!\mathit{olarized}}
\def\Complexesstack{\mathcal{C}\!\mathit{omplexes}}
% \Pic is the operator that assigns to X its picard group, usage \Pic(X)
% \Picardstack_{X/B} denotes the Picard stack of X over B
% \Picardfunctor_{X/B} denotes the Picard functor of X over B
\def\Pic{\mathop{\mathrm{Pic}}\nolimits}
\def\Picardstack{\mathcal{P}\!\mathit{ic}}
\def\Picardfunctor{\mathrm{Pic}}
\def\Deformationcategory{\mathcal{D}\!\mathit{ef}}

% --- appended 2026-10-12 (Chief): the project own definitions of the reference/slogan/history
% display environments, copied mechanically from the shared Stacks preamble lines 8-12
% (\usepackage{verbatim} plus the three \newenvironment lines, which the chapter preamble keeps
% verbatim from the project preamble). Nothing added, removed or redefined: these environments
% carry no mathematical content and the project itself typesets them as comments. Without these
% three lines a counted statement whose own text uses them fails with Environment slogan undefined.
\usepackage{verbatim}
\newenvironment{reference}{\comment}{\endcomment}
\newenvironment{slogan}{\comment}{\endcomment}
\newenvironment{history}{\comment}{\endcomment}
\usepackage{cleveref}
\crefname{originallemma}{original lemma}{original lemmas}
\Crefname{originallemma}{Original lemma}{Original lemmas}
\makeatletter
\newlabel{more-algebra-lemma-J-2}{{07PC}{source}{}{sourcepointer-1}{}}
\newlabel{more-algebra-lemma-J-2@cref}{{[originallemma][0][]07PC}{[1][0][]source}{}{}{}}
\makeatother
\crefname{originallemma}{original lemma}{original lemmas}
\Crefname{originallemma}{Original lemma}{Original lemmas}
\makeatletter
\newlabel{lemma-intersection-regular-with-closed}{{07P9}{source}{}{sourcepointer-2}{}}
\newlabel{lemma-intersection-regular-with-closed@cref}{{[originallemma][0][]07P9}{[1][0][]source}{}{}{}}
\makeatother
\crefname{originallemma}{original lemma}{original lemmas}
\Crefname{originallemma}{Original lemma}{Original lemmas}
\makeatletter
\newlabel{algebra-lemma-make-separably-generated}{{04KM}{source}{}{sourcepointer-3}{}}
\newlabel{algebra-lemma-make-separably-generated@cref}{{[originallemma][0][]04KM}{[1][0][]source}{}{}{}}
\makeatother
\crefname{originallemma}{original lemma}{original lemmas}
\Crefname{originallemma}{Original lemma}{Original lemmas}
\makeatletter
\newlabel{lemma-J-0-goes-up}{{07PB}{source}{}{sourcepointer-4}{}}
\newlabel{lemma-J-0-goes-up@cref}{{[originallemma][0][]07PB}{[1][0][]source}{}{}{}}
\makeatother
\crefname{originallemma}{original lemma}{original lemmas}
\Crefname{originallemma}{Original lemma}{Original lemmas}
\makeatletter
\newlabel{lemma-J-0-goes-down}{{07PA}{source}{}{sourcepointer-5}{}}
\newlabel{lemma-J-0-goes-down@cref}{{[originallemma][0][]07PA}{[1][0][]source}{}{}{}}
\makeatother
\title{Four equivalent formulations of the J-2 property}
\author{The Stacks Project Authors}
\date{stacks-project:more-algebra.tex:local-snapshot-20261012; GNU FDL 1.2 or later}
\begin{document}
\maketitle
\setcounter{subsection}{1}
\noindent\textit{Editorial note.} All delivered mathematical text is the author's own, copied from the pinned Stacks Project chapter file more-algebra.tex; no statement or proof step was rewritten, reordered or removed.\par\medskip
\noindent\textit{Editorial note.} The chapter file pinned-original-source.tex carries only \textbackslash{}input\{preamble\} on its first line, so no preamble text lives inside the chapter. The retained file sources/stacks-more-algebra/retained-preamble.tex is a mechanical, line-for-line copy of the project's single shared preamble as staged in the snapshot (Xy-pic setup plus theorem declarations and mathematical macros). Nothing was added, removed or redefined, and the excerpt serves only as preamble, never as a span coordinate.\par\medskip
\noindent\textit{Editorial note.} Both counted spans start at the \textbackslash{}begin line of their environment, so the statement and the author's proof are delivered whole; surrounding prose of the section is deliberately outside the counted spans and is carried as standard prerequisites instead.\par\medskip
\noindent\textit{Editorial note.} The citation fields (source\_id, source\_version, authors, source\_url, license.name, source\_title) are identical to those of every other card this lane has produced from this chapter, so the batch selector treats them as one source.\par\medskip
\noindent\textit{Editorial note.} Benign-screen note: the proof contains the phrase 'it is clear that' / 'is obvious' once, used as a legitimate step inside a complete argument; the author proves the statement, so this is not a hint-only proof.\par\medskip
\noindent\textit{Allowed standard prerequisites (editorial).}\begin{itemize}
\item A J-1 domain is J-0, and the criterion for J-2 through quotients by primes. Lemma \textbackslash{}ref\{lemma-intersection-regular-with-closed\}, cited at proof lines 11710-11712. Give two of the implications.
\item Existence of a diagram of finite purely inseparable extensions making the resulting extension separable. Algebra, Lemma \textbackslash{}ref\{algebra-lemma-make-separably-generated\}, cited at proof lines 11722-11724. Produces the extension used in the final step.
\item Ascent and descent of the J-0 property. Lemma \textbackslash{}ref\{lemma-J-0-goes-up\} and Lemma \textbackslash{}ref\{lemma-J-0-goes-down\}, cited at proof lines 11731 and 11732-11734. Transfer J-0 between the auxiliary domain and the original one.
\end{itemize}\medskip
\ifdefined\hypertarget\hypertarget{sourcepointer-1}{}\fi
\noindent\textit{Source reference.} Original lemma 07PC: focused-seeds/AGP/sources/stacks-snapshot-20261011/tags.txt, source lines 9792--9792. The lemma's own canonical tag, quoted so that readers can resolve the label inside the delivered text; it points at the counted statement itself and carries no proof content.\par\medskip
\ifdefined\hypertarget\hypertarget{sourcepointer-2}{}\fi
\noindent\textit{Source reference.} Original lemma 07P9: focused-seeds/AGP/sources/stacks-snapshot-20261011/tags.txt, source lines 9789--9789. Navigation pointer to a statement the counted proof refers to but does not contain; nothing load-bearing is delegated to it.\par\medskip
\ifdefined\hypertarget\hypertarget{sourcepointer-3}{}\fi
\noindent\textit{Source reference.} Original lemma 04KM: focused-seeds/AGP/sources/stacks-snapshot-20261011/tags.txt, source lines 5978--5978. Navigation pointer to a statement the counted proof refers to but does not contain; nothing load-bearing is delegated to it.\par\medskip
\ifdefined\hypertarget\hypertarget{sourcepointer-4}{}\fi
\noindent\textit{Source reference.} Original lemma 07PB: focused-seeds/AGP/sources/stacks-snapshot-20261011/tags.txt, source lines 9791--9791. Navigation pointer to a statement the counted proof refers to but does not contain; nothing load-bearing is delegated to it.\par\medskip
\ifdefined\hypertarget\hypertarget{sourcepointer-5}{}\fi
\noindent\textit{Source reference.} Original lemma 07PA: focused-seeds/AGP/sources/stacks-snapshot-20261011/tags.txt, source lines 9790--9790. Navigation pointer to a statement the counted proof refers to but does not contain; nothing load-bearing is delegated to it.\par\medskip
\begin{lemma}
\label{lemma-J-2}
Let $R$ be a Noetherian ring. The following are equivalent
\begin{enumerate}
\item $R$ is J-2,
\item every finite type $R$-algebra which is a domain is J-0,
\item every finite $R$-algebra is J-1,
\item for every prime $\mathfrak p$ and every finite purely inseparable
extension $L/\kappa(\mathfrak p)$ there exists a finite
$R$-algebra $R'$ which is a domain, which is J-0, and whose field
of fractions is $L$.
\end{enumerate}
\end{lemma}

\begin{proof}
It is clear that we have the implications (1) $\Rightarrow$ (2) and
(2) $\Rightarrow$ (4). Recall that a domain which is
J-1 is J-0. Hence we also have the implications
(1) $\Rightarrow$ (3) and (3) $\Rightarrow$ (4).

\medskip\noindent
Let $R \to S$ be a finite type ring map and let's try to show $S$ is J-1. By
Lemma \ref{lemma-intersection-regular-with-closed} it suffices
to prove that $S/\mathfrak q$ is J-0 for every prime $\mathfrak q$
of $S$. In this way we see (2) $\Rightarrow$ (1).

\medskip\noindent
Assume (4). We will show that (2) holds which will finish the proof.
Let $R \to S$ be a finite type ring map with $S$ a domain.
Let $\mathfrak p = \Ker(R \to S)$. Let $K$ be the fraction field of $S$.
There exists a diagram of fields
$$
\xymatrix{
K \ar[r] & K' \\
\kappa(\mathfrak p) \ar[u] \ar[r] & L \ar[u]
}
$$
where the horizontal arrows are finite purely inseparable field extensions
and where $K'/L$ is separable, see
Algebra, Lemma \ref{algebra-lemma-make-separably-generated}.
Choose $R' \subset L$ as in (4) and let
$S'$ be the image of the map $S \otimes_R R' \to K'$.
Then $S'$ is a domain whose fraction field is $K'$, hence
$S'$ is J-0 by Lemma \ref{lemma-J-0-goes-up} and our choice of $R'$.
Then we apply Lemma \ref{lemma-J-0-goes-down} to see that $S$
is J-0 as desired.
\end{proof}

\end{document}
